\begin{claim}{C.34}{Proved}{The median contraction on all nonopen orbits}
Let $f_t$ be a based null-homotopic loop in
$\operatorname{Homeo}^+(\R^2)$, and use $P,G,Y$ from
C.28 and C.31. The orbitwise bijections $Q_{s,t}:Y\to Y$
defined in C.31 are homeomorphisms. Both maps
\[
 [0,1]^2\times Y\longrightarrow Y,\qquad
                  (s,t,x)\longmapsto Q_{s,t}^{\pm1}(x)
                                                        \tag{C.34.1}
\]
are continuous. Thus these formulas give a based
compact-open null-homotopy on the entire closed nonopen
locus, commuting with every homeomorphism commuting
with all $f_t$.
In particular, a hyperbolic centralizer loop whose
canonically lifted generated group has no open orbit
is null-homotopic in the original surface centralizer.
No extension over the omitted open orbits is asserted.
\end{claim}
\begin{proof}
We prove continuity at a nonfixed orbit. Continuity at
the common fixed set, bijectivity on $Y$, equivariance,
and the exact inverse identity have already been proved
in C.31.

\emph{Short paths moving in either direction.}
Let $q$ be a point of a nonopen nontrivial orbit. It lies
on an intrinsic line or a circle by C.28. There are
paths $g_s$ from the identity, formed from two factors
in $P$, which move $q$ a nonzero amount in either
chosen intrinsic direction while keeping its whole
track inside any prescribed sufficiently small
intrinsic arc about $q$.
To see this, $f_t(q)$ is not identically $q$; otherwise
every generator, and hence every element of $G$, would
fix $q$. Choose a component $(a,b)$ of
$\{t:f_t(q)\ne q\}$ and a sufficiently short initial
subinterval $[a,a+\eta]$. The path
\[
          h_s=f_{a+s\eta}f_a^{-1},\qquad 0\leq s\leq1,
                                                        \tag{C.34.2}
\]
starts at the identity, and $h_s(q)=f_{a+s\eta}(q)$.
C.28 gives compact-open continuity on an intrinsic
line; the circle restriction is continuous as well.
Consequently $h_s$ and $h_s^{-1}$ can be made uniformly
close to the identity on an intrinsic neighborhood of
$q$. Their endpoint displacements at $q$ have opposite
intrinsic signs. One of these two paths has the chosen
sign, and its track stays in the prescribed arc.
The paths also become uniformly small on each fixed
ambient compact set as $\eta\to0$.
Every such $g_s$ is a product of two elements of $P$.
For any $z\in K_n(x)=P^nx$,
\[
                        g_s(z)\in K_{n+2}(x).
                                                        \tag{C.34.3}
\]
In particular, the same compact-word argument as in
C.28 proves intrinsic continuity of its track on
every nonopen orbit, including when $z$ varies along
a compact intrinsic subarc.

\emph{A small cap cannot persist at a moving limit point.}
We give the local argument used below. Straighten a
compact arc of a limit orbit, extending it a little
at both ends, so that it is an interval of the
horizontal axis in a plane chart. Suppose compact
intrinsic arcs $A_j$ in nearby nonopen orbits lie in
strips of height tending to zero around this interval.
Let $q_j\in A_j$ attain the minimum $m_j$ of the
horizontal coordinate on $A_j$, suppose
$q_j\to q=(m,0)$, and suppose both endpoints of
$A_j$ have horizontal coordinates at least $m+\delta$
for a fixed $\delta>0$.
Take a rectangle
\[
          U=(m-\rho,m+\rho)\times(-\varepsilon,\varepsilon),
          \qquad 0<\rho<\delta/3,
                                                        \tag{C.34.4}
\]
whose closure stays in the chart. For large $j$,
$q_j$ lies in $U$, all of $A_j$ has vertical coordinate
of absolute value less than $\varepsilon/2$, and its
horizontal minimum is greater than $m-\rho/2$.
Starting at $q_j$ in each intrinsic direction towards
the two endpoints of $A_j$, take the first point
outside $U$. Both first-exit points lie on its right
edge. They are distinct points of the embedded arc.
Let $J_j$ be the closed subarc between these points.
Its interior lies in $U$, and every point of $J_j$
has horizontal coordinate at least $m_j$.

Use (C.34.2), taking the inverse when needed, to choose
a path $g_s$ whose $q$-track stays in $U$ and whose
endpoint is strictly to the left of $q$. This choice
is possible because $q$ is nonfixed. The track is a
compact subset of $U$, so joint continuity gives
$g_s(q_j)\in U$ for every $s$ and every sufficiently
large $j$.
Its intrinsic continuity from (C.34.3) now matters:
a continuous path in the intrinsic line or circle
starting at an interior point of $J_j$ cannot leave
$J_j$ without meeting one of its two endpoints.
Both endpoints lie outside the open rectangle $U$.
Thus the entire $g_s(q_j)$-track stays in $J_j$.
At $s=1$ its horizontal coordinate is at least $m_j$,
whereas
\[
        g_1(q_j)\longrightarrow g_1(q),\qquad
        \operatorname{pr}_1(g_1(q))<m=\lim_jm_j.
                                                        \tag{C.34.5}
\]
This contradiction excludes the cap.
The argument with horizontal maxima and a path moving
to the right is identical. The two first-exit points
are the explicit obstruction to escaping the cap
along another part of the same orbit.

The same argument excludes a whole Jordan-curve orbit
contained in such a shrinking strip and having a
nondegenerate interval as its limit: use a global
minimum of the horizontal coordinate on the circle,
and in its two directions take first exits towards
a point whose coordinate is bounded above that
minimum. The resulting subarc $J_j$ again has
two distinct endpoints outside $U$.

\emph{No excess in subarcs near an intrinsic line.}
Fix a point $p$ on an intrinsic line orbit and
take $x_j\to p$ in $Y$. For each fixed $n$, compactness
of $P^n$ and continuity of word evaluation give
\[
                         K_n(x_j)\longrightarrow K_n(p)
              \quad\hbox{in Hausdorff distance}.         \tag{C.34.6}
\]
The limit is a nondegenerate compact arc. Straighten
a slightly larger compact arc of the limit orbit.
Then (C.34.6) places the nearby $K_n(x_j)$ in shrinking
horizontal strips in that chart. Each $K_n(x_j)$ is
an arc, a point, or a whole circle, by C.28 and
connectedness. The whole-circle possibility is
excluded for all sufficiently large $j$ by the
preceding cap argument.

Suppose $u_j,v_j\in K_n(x_j)$ tend to $u,v\in K_n(p)$.
Write $[u_j,v_j]_j$ for the subarc contained in
$K_n(x_j)$. Then
\[
 \begin{gathered}
 z_j\in[u_j,v_j]_j,\quad z_j\longrightarrow z
             \quad\Longrightarrow\quad z\in[u,v].
                                                        \tag{C.34.7}
 \end{gathered}
\]
Indeed, an alleged limit point strictly to the left
of both $u,v$ produces a horizontal minimum of
$[u_j,v_j]_j$ whose two endpoints stay a fixed
distance to its right. A subsequence of these
minimum points converges, by (C.34.6), to a point
of the nonfixed limit orbit. This is the forbidden
cap. A point to the right is excluded by the
maximum version. The assertion includes $u=v$:
then all of the intervening subarc tends to $u$.
This proves the needed continuity of subarcs using
the path action; Hausdorff convergence alone has
not been used to infer it.

It follows that medians of three points in $K_n(x_j)$
converge to the median of their limits. Each median
lies in all three pairwise connecting subarcs.
Every subsequential limit therefore lies, by
(C.34.7), in all three corresponding limit intervals.
Their intersection is the single limit median,
including when two or three of the limiting points
coincide.

\emph{A circle limit: the precise local covering.}
Let the orbit of $p$ be a Jordan curve $C$. Choose a
collar annulus $W$ of $C$, with coordinates
$(\R/\Z)\times(-1,1)$ sending $C$ to the middle circle,
and write
\[
 \pi_W:\widetilde W=\R\times(-1,1)\longrightarrow W
                                                        \tag{C.34.8}
\]
for its universal covering. The compact family
$f_t$ preserves $C$. Therefore a sufficiently small
open disc $V$ about $p$ satisfies
$f_t(V)\subset W$ for all $t$.
The actual map
\[
 [0,1]\times V\longrightarrow W,\qquad
                         (t,x)\longmapsto f_t(x)
                                                        \tag{C.34.9}
\]
has a unique continuous lift $\widehat f_t(x)$ after
choosing a lift of the initial inclusion $V\to W$.
Its endpoint differs from its starting point by
an integer deck translation depending continuously
on $x\in V$. It is zero at $p$, by C.29's exact
plane-to-circle winding comparison, so it is zero
on all of the connected disc $V$.
Thus these are closed lifted point tracks.
For $x_j\to p$ in $V\cap Y$ put
\[
 T_j(r)=\{\widehat f_u(x_j):0\leq u\leq r\},\qquad
 T(r)=\{\widehat f_u(p):0\leq u\leq r\}.
                                                        \tag{C.34.10}
\]
Joint continuity on the compact parameter interval
shows $T_j(r_j)\to T(r)$ in Hausdorff distance
whenever $r_j\to r$. The set $T(1)$ is a
nondegenerate compact interval on the line
$\R\times\{0\}$.

We verify the intrinsic meaning of the lifted
nearby tracks. An intrinsic line orbit, or a
proper component of a circle orbit inside $W$,
has an intrinsic interval as its component through
$x_j$. Its continuous point track stays in that
component and lifts to an injectively parametrized
interval in $\widetilde W$, with embedded compact
subintervals. If the entire circle orbit lies in
$W$, it is a Jordan curve. Its winding in $W$
is $0$ or $\pm1$: the planar Jordan winding
around the complementary inner disc has exactly
these values. Its lifted components are therefore
Jordan curves in the first case and intrinsic
lines in the second.
Consequently $T_j(1)$ is a compact arc, a point,
or a whole lifted Jordan curve, with the
corresponding intrinsic topology.

The cap argument applies in this covering chart
with the original group paths. Here is the exact
local map needed to justify that use. At any
limiting extremum $q\in T(1)$ choose the small
rectangle $U$ of (C.34.4) with horizontal width
less than one, so that $\pi_W|_U$ is injective.
Choose a path (C.34.2) at $\pi_W(q)$ with the
required intrinsic sign and sufficiently small
that its images of a compact neighborhood of
$\pi_W(q)$ stay in $W$.
The map
\[
 (s,z)\longmapsto g_s(\pi_W(z))
                                                        \tag{C.34.11}
\]
on that small lifted neighborhood has a unique
continuous lift starting with $z\mapsto z$.
Its $q$-track lies in $U$ and has the required
horizontal displacement. For $q_j\to q$ its
whole $q_j$-track also stays in $U$, by compact
parameter continuity. Downstairs it is an
intrinsically continuous path in the original
orbit, by (C.34.3). Its lift cannot leave the
cap subarc without meeting one of its first-exit
endpoints on $\partial U$. Thus (C.34.5) applies
unchanged. No action of $G$ on the entire
annulus or its covering has been assumed.

This excludes a whole lifted Jordan curve for
large $j$ and proves (C.34.7), and hence median
continuity, for compact subarcs of $T_j(1)$.
These intrinsic interval operations are exactly
the ones used by C.27 and C.29. For a line or
a proper circle component the lift of its
interval is order preserving or reversing.
For an essential circle the annular covering
is its universal line covering. For a
contractible circle, the now proper arc
$T_j(1)$ lies in a single interval of its
universal line covering; a path leaving that
interval would have to cover the whole circle.
In all cases the based circle lift or intrinsic
line formula, followed by projection through
(C.34.8), gives the original $Q_{s,t}(x_j)$.
Median invariance under reversal removes any
choice of the interval orientation.

\emph{Continuity of the prefix-clamping formula.}
We can now handle a sequence
$(s_j,t_j,x_j)\to(s,t,p)$ with $p$ nonfixed.
For a line limit use the actual point tracks
$T_j(r)=\{f_u(x_j):0\leq u\leq r\}$ inside
$K_1(x_j)$. For a circle limit use (C.34.10).
In either case $T_j(1)$ is eventually an
intrinsic compact arc, its subarcs satisfy
(C.34.7), and its prefix tracks have their
stated Hausdorff limits.
Set $r_j=1-s_j$ and $r=1-s$.
Let $a_j,b_j$ be the two unordered endpoints
of the prefix arc $T_j(r_j)$, allowing equality.
Every subsequence has a further subsequence
on which both endpoints converge, say to
$a,b\in T(r)$. By (C.34.7), the Hausdorff
limit of their intervening arc is contained
in $[a,b]$. It is also exactly $T(r)$ by
the prefix-track continuity. Since $a,b$
belong to $T(r)$, their interval is contained
in $T(r)$. Therefore
\[
                         T(r)=[a,b].                    \tag{C.34.12}
\]
In particular the limiting unordered pair is
precisely the endpoint pair of the prefix
interval, including a degenerate prefix.

The formula of C.27 or its C.29 lift is
the median of $a_j,b_j$ and the current
track point at $t_j$. The median continuity
just proved and (C.34.12) give its limit as
the median of the limiting prefix endpoints
and the current point at $t$. This is
$Q_{s,t}(p)$, with projection through
(C.34.8) in the circle case.
Every subsequence has this same limit.
The domains are metrizable, so this proves
joint continuity on $[0,1]^2$ times the
nonfixed part of $Y$. C.31 supplies the
remaining fixed points.

Apply the same argument to the based inverse
loop $u\mapsto f_u^{-1}$. It is ambient-null,
generates the same group $G$, and has the
same closed locus $Y$. The exact identity
$Q[f]^{-1}=Q[f^{-1}]$ in C.31 proves joint
continuity of the inverse family in
(C.34.1). Thus all the orbitwise bijections
are homeomorphisms of $Y$. The compact-open
continuity follows from the evaluation
exponential law, since $Y$ is a closed,
locally compact subspace of the plane.
The based endpoints and equivariance are
those already proved orbitwise in C.31.

Finally, when the canonical lifted group of
a hyperbolic surface loop has no open orbit,
$Y=\widetilde U\cong\R^2$.
The obtained plane homeomorphisms commute
with the full operator lift group, including
all deck transformations. They descend by
the exact equations
\[
                  \overline Q_{s,t}\,p_U=p_UQ_{s,t}
                 \quad(p_U:\widetilde U\to U)
                                                        \tag{C.34.13}
\]
to homeomorphisms of $U$ commuting with every
original operator. Their joint continuity
and that of their inverses descend through
the open covering quotient, giving
compact-open continuity downstairs.
The endpoints in (C.34.13) are the original
loop and the identity, with both based
edges fixed. Hence this is its null-homotopy
inside the original centralizer.
At every parameter it lies in the identity
component by these same paths.
Equivalently, its lift is the canonical
one: two deck-commuting lifts of such a
map differ by a central deck transformation,
and the hyperbolic surface group has trivial
center. This proves the stated corollary
with the actual lift and descent maps.
\end{proof}

\paragraph{Remaining global extension.}
C.34 establishes the full nonopen-locus
homeomorphism family and its inverse.
C.33 provides intrinsic equivariant
contractions on each open orbit.
C.35 proves that this particular median
family need not extend equivariantly over
an adjacent open orbit. C.36 and C.37
construct compatible alternatives on the
closed union $X_2$ of open orbits with
two-dimensional proper operator image,
including every infinite accumulation.
Matching that boundary family to the
exterior of $X_2$, and treating the other
proper operator images at their frontiers,
remain the assembly requirements in
C.12, whose general status remains open.
